\addcontentsline{toc}{section}{CAPÍTULO IX}
\chapter{MÉTRICA PRINCIPAL DE VALIDACIÓN}

\section{Definición}

La métrica principal del MVP es la \textbf{tasa de \textit{fast-track}:} el porcentaje de liquidaciones hospitalarias que el sistema evalúa, clasifica como de bajo riesgo y presenta al auditor con una recomendación fundamentada que este puede validar en aproximadamente 1 minuto. La meta es alcanzar una tasa del 50–60 \%. El rango de 50–60 \% representa la banda de operación óptima estimada. Una tasa inferior al 50 \% indicaría un impacto operativo insuficiente. Una tasa superior al 60 \% pero inferior al 75 \% es aceptable si no se observa un incremento en la tasa de falsos negativos, ya que implica que el sistema está liberando más capacidad del auditor de la inicialmente prevista.

Esta métrica refleja directamente la propuesta de valor de \textbf{\textit{NexusCare}}. Una tasa de \textit{fast-track} del 55 \% significa que, de un lote de 101 liquidaciones, aproximadamente 56 se resuelven en alrededor de 1 minuto cada una, lo que libera al auditor para concentrarse en los casos restantes que requieren un análisis profundo. El impacto operativo se mide a nivel de lote porque así opera el auditor: el valor se materializa cuando un lote completo se cierra más rápido, no cuando un caso aislado se procesa en menos tiempo.

La línea base se estableció durante Capstone 1: un lote de 101 liquidaciones hospitalarias requiere 75.8 horas de trabajo manual, aproximadamente 45 minutos por caso. La cifra de 45 minutos proviene de entrevistas con auditores médicos de MediVida durante Capstone 1, quienes reportaron un rango de 30 a 45 minutos para las liquidaciones hospitalarias, según la complejidad del caso. Se adoptó el extremo superior como supuesto conservador; no se midió mediante un estudio formal de tiempo y movimiento. Con una tasa de \textit{fast-track} del 55–60 \%, el tiempo total por lote se reduce a un rango estimado de 28 a 35 horas, equivalente a una reducción del 40–60 \%, según el tiempo efectivo que el auditor dedique a los casos de revisión rápida y a la auditoría completa.

\section{Métricas secundarias}

 \textbf{Tasa de falsos negativos (umbral máximo: 2 \%)} Mide el porcentaje de liquidaciones de riesgo medio o alto que el sistema clasifica incorrectamente como de bajo riesgo y recomienda para \textit{fasttrack}. A diferencia de los falsos positivos, que generan ineficiencia operativa, un falso negativo implica que una liquidación con inconsistencias clínicas o indicios de fraude pasa con una revisión mínima del auditor. Es la métrica de seguridad más crítica del sistema. El umbral del 2 \% se define como requisito no negociable: por encima de ese nivel, el riesgo regulatorio y reputacional para la aseguradora supera el beneficio operativo. Su medición precisa requiere un piloto con auditor real; en el alcance del MVP se estima indirectamente a partir del \textit{Golden Set}.

\textbf{Tasa de falsos positivos} Corresponde a liquidaciones clasificadas incorrectamente como de alto riesgo que, tras revisión del auditor, resultan ser casos rutinarios. Los falsos positivos tienen un costo operativo directo y un costo de confianza. Esta métrica solo puede medirse con precisión en un piloto con auditor real; en el alcance del MVP se estima a partir del \textit{Golden Set}.

\textbf{Concordancia auditor-sistema} Es el porcentaje de veces que el auditor acepta la recomendación del sistema sin modificarla. Valida el supuesto crítico de adopción y confianza. La interfaz React registra las acciones del auditor para calcular esta métrica en producción.

\textbf{Tiempo total por lote} Es un KPI derivado de la tasa de \textit{fast-track} y del tiempo de revisión de los casos no \textit{fast-track}. Se estima multiplicando el número de casos en cada categoría por el tiempo de revisión correspondiente. La meta es una reducción de al menos el 40 \% respecto a la línea base de 75.8 horas. Esta meta es conservadora: asume que todos los casos de auditoría completa siguen tomando 45 minutos y que los de revisión rápida toman 10 minutos. Si el auditor desarrolla confianza en el sistema y reduce el tiempo de los casos de revisión rápida, la reducción real podría superar el 60 \%.

\section{Método de medición}

La validación del MVP se basa en simulación y proyección, no en medición directa con un auditor en producción. Esta limitación se declara explícitamente.

La simulación no valida que los puntajes del LLM caerán en los rangos esperados; eso solo puede confirmarlo el \textit{benchmark} con el modelo real y, en última instancia, un piloto con auditor. Lo que la simulación sí valida es la lógica del motor de decisión: dado un puntaje X y un resultado de Fase 1 Y, el sistema debe producir la categoría correcta. Los puntajes simulados (\textit{score} 20 para cobertura completa, \textit{score} 40 para condicionales, \textit{score} 55 para cobertura parcial) se derivan de la distribución observada en el benchmark de Sonnet, no de un supuesto arbitrario. Sin embargo, la distancia entre la distribución del benchmark y la distribución real en operación es una incógnita que solo un piloto puede resolver.

La tasa de \textit{fast-track} se proyecta mediante el \textit{test} test\_projected\_rate\_meets\_target, que ejecuta Fase 1 sobre las 101 liquidaciones con los planes MediVida Más 65 y MediVida Premium, simula puntajes de Fase 2 según el resultado de cobertura y verifica que la tasa resultante se encuentre entre 55 \% y 75 \%. El \textit{test} pasa con las configuraciones actuales, proyectando una tasa en el rango del 60–68 \%. El techo del 75 \% existe precisamente para detectar calibración insuficiente: una tasa excesivamente alta sugiere que el sistema no está siendo lo suficientemente riguroso en la evaluación de riesgo.

El tiempo total por lote se estima a partir de la distribución proyectada de decisiones, multiplicando el número de casos en cada categoría por el tiempo estimado de revisión: 1 minuto para FAST\_TRACK, 10 minutos para REVISION\_RAPIDA, 45 minutos para AUDITORIA\_COMPLETA y 2 minutos para RECHAZO\_ADMINISTRATIVO. La validación real requeriría un piloto controlado donde un grupo de auditores procese un lote con \textbf{\textit{NexusCare}} y otro sin él, midiendo el tiempo total de cierre del lote. Ese piloto queda fuera del alcance del MVP académico.
